%% Draft v2 (ICML 2025 style): Approach, Evaluation, Ablations only. Rewritten around the information ladder in
%% ../UPDATED PLAN.md. paper.tex (the older draft) is not edited. Abstract, Introduction, Related Work, Conclusion come last.
%% Red text marks numbers that are still running; `grep -n "pend\|todo\|tbd" paper_v2.tex` must be empty before submitting.
\documentclass{article}
\usepackage{microtype}
\usepackage{graphicx}
\usepackage{booktabs}
\usepackage{multirow}
\usepackage{xcolor}
\usepackage{hyperref}
\newcommand{\theHalgorithm}{\arabic{algorithm}}
\usepackage[accepted]{icml2025}   % non-anonymous, as the course allows
\usepackage{amsmath,amssymb}
\usepackage[capitalize,noabbrev]{cleveref}

\newcommand{\todo}[1]{\textcolor{red}{[TODO: #1]}}
\newcommand{\tbd}{\textcolor{red}{??}}
\newcommand{\pend}{\textcolor{red}{pending}}
\newcommand{\fig}[2]{\IfFileExists{#1}{\includegraphics[width=#2]{#1}}%
  {\fbox{\parbox[c][3cm][c]{#2}{\centering\small\todo{\detokenize{#1}}}}}}

\icmltitlerunning{Pixels or Keypoints? An Information Ladder for Low-Data Sign Recognition}

\begin{document}
\twocolumn[
\icmltitle{Pixels or Keypoints? An Information Ladder\\for Low-Data Isolated Sign Recognition}
\begin{icmlauthorlist}
\icmlauthor{Narelle Ho}{ucb}
\icmlauthor{Cheng Yao Lee}{ucb}
\end{icmlauthorlist}
\icmlaffiliation{ucb}{CS289A, University of California, Berkeley}
\icmlcorrespondingauthor{Cheng Yao Lee}{\todo{email}}
\icmlkeywords{Sign language recognition, WLASL, pose estimation, video classification}
\vskip 0.3in
]
\printAffiliationsAndNotice{}

\begin{abstract}
\todo{written last}
\end{abstract}

%% Introduction and Related Work: written after the results are final.

%=====================================================================
\section{Approach}\label{sec:approach}
%=====================================================================
We ask whether, with only a few training videos per sign, a recognizer should learn from \emph{pixels} or
from a \emph{skeleton}: keypoints found by a pretrained hand detector. Published comparisons change the data,
the preprocessing and the model at the same time, so a score difference does not say where it comes from.
We therefore build an \textbf{information ladder}. Every step changes one thing, the same videos and the same
16 frames per video are used throughout, and the step sizes show what matters (\cref{fig:ladder}).

\begin{figure*}[t]
\centering
\fig{figs/ladder_diagram.png}{0.86\textwidth}
\caption{The information ladder. Orange: the network reads pixels (ResNet-18 features of each frame, then a small
transformer). Blue: the same small transformer reads keypoints. Arrows name what changes: where the network looks,
whether it knows where the hand is, whether the hand is given as pixels or as keypoints, and whether the upper body is
included. Labels S3, S5, X1, S1, X2 are used in the tables.}
\label{fig:ladder}
\end{figure*}

\subsection{Task and data}
We use WLASL100 \citep{li2020wlasl}: isolated signs, 100 words. Only 1{,}013 of the 2{,}038 videos could be
obtained, so all results use the official split restricted to them: 748 training, 165 validation and 100 test videos
(7.5 training videos per sign on average). The median video is 656$\times$400 pixels and 65 frames long. The test set
covers 72 of the 100 signs, and 97 of its 100 videos are by signers who also appear in training
(\cref{sec:limits}). Choices (the best training epoch) are made on the validation videos only; test videos are used
for reporting.

\subsection{Shared input}\label{sec:shared}
Each video is reduced once to \textbf{16 evenly spaced frames}, resized to 224$\times$224 with black bars so the
shape is kept. Both groups use exactly these frames.
\begin{itemize}
\item \textbf{Pixels.} The 224$\times$224 frame, or two 112$\times$112 crops per frame, one per hand, cut from it.
      A crop is a square of 1.6 times the hand's bounding box (at least 32 pixels), enlarged to 112$\times$112.
      The crops are raw pixels; nothing is drawn on them.
\item \textbf{Keypoints.} MediaPipe Hands \citep{lugaresi2019mediapipe,zhang2020mediapipehands} is run on each of the 16
      frames separately, with no tracking and no neighbouring frames, so the keypoints cannot contain anything the
      pixel model cannot see. Each hand has 21 points (x, y, depth). In 23\% of the frames both hands are found,
      in 41\% one hand, in 36\% none.
      For the upper body (step X2) MediaPipe Pose \citep{bazarevsky2020blazepose} is run the same way; we keep nine
      points (nose, shoulders, elbows, wrists, hips). A person is found in 98.5\% of frames, and in at least some frame
      of every video.
\end{itemize}
The crops use the Hands detector to decide \emph{where} to cut, so the crop models are pose-guided pipelines too;
their network sees only pixels.

\subsection{Keypoint models}
Each frame becomes a vector of 132 numbers: for each hand, the 21 points relative to the wrist and divided by the palm
length (126), a ``hand found'' flag per hand (2), and the position of each wrist in the frame (4). A hand that is
missing is zero with its flag set. In step X2 we append the nine body points relative to the shoulder midpoint and
divided by shoulder width (18 numbers) and a ``body found'' flag. Features are standardised with training statistics.
The \textbf{small transformer} \citep{vaswani2017attention} has one layer, width 128, 4 heads, feed-forward width 256,
dropout 0.3, learned position embeddings over the 16 frames, mean over frames, and a linear layer to 100 classes
(165k parameters). It is trained from scratch.

\subsection{Pixel models}
\textbf{ResNet-18} \citep{he2016deep} pretrained on ImageNet \citep{deng2009imagenet} turns each frame, or each hand crop
(the same weights for both hands), into 512 numbers. The 16 vectors, plus the hand-found flags where used, go through
the \emph{same} transformer as above (one layer, width 128). The whole network is fine-tuned on our 748 videos
(11.4 to 11.5M parameters). In step X1 the crop model also receives the same four wrist-position numbers as the keypoint
model, standardised, zeroed when the hand is missing, and given the same random rotation, zoom and shift as the keypoints.
\textbf{I3D} \citep{carreira2017quo} is the network used for the WLASL video baseline \citep{li2020wlasl}, pretrained on
Kinetics \citep{kay2017kinetics}; we fine-tune it on the 16 whole frames (12.4M parameters).

\subsection{Training and statistics}
All models: label smoothing 0.1 \citep{szegedy2016rethinking}, AdamW \citep{loshchilov2019decoupled} with weight decay
0.01, gradient clipping 1.0, the best epoch chosen on validation accuracy, three seeds (0, 1, 2).
Keypoint models: learning rate $10^{-3}$, batch 64, up to 200 epochs with patience 40. Fine-tuned pixel models: learning
rate $10^{-4}$ for the pretrained network and $10^{-3}$ for the transformer and new layers, cosine schedule, mixed
precision, batch 8, up to 40 epochs with patience 10. Augmentation is one random transform per video: up to 10\% of
frames trimmed from each end, rotation $\pm15^\circ$, zoom $\pm10\%$, shift $\pm5\%$, small noise; keypoints get the same
transform in coordinate form. There is no horizontal flip, which would swap the signing hand. \emph{No hyper-parameter
search}: one recipe per group, so differences between rows are not tuning differences within a group, and tuning
across groups is a stated limitation (\cref{sec:limits}).

We report the number of the 100 test videos classified correctly (``accuracy''), the share with the right sign among the
five best guesses (top-5), and the spread over seeds. The seed spread ignores the uncertainty from having only 100 test
videos (about $\pm$10 points for one accuracy), so \textbf{differences between models are given with a paired 95\% interval}:
the score of a video is the share of seeds that got it right, the difference is averaged over the same 100 videos, and the
interval is a bootstrap over videos with 10{,}000 draws \citep{efron1994introduction}. Three comparisons were fixed
before the new runs finished (``primary''): keypoints vs.\ crops with the same location (S1 vs.\ X1), the value of location
for the crop model (X1 vs.\ S5), and the value of the upper body for the keypoint model (X2 vs.\ S1). All others are
exploratory.

%=====================================================================
\section{Evaluation}\label{sec:eval}
%=====================================================================
\subsection{Main comparison}
\begin{table*}[t]
\centering\small\setlength{\tabcolsep}{4pt}
\begin{tabular}{llcccrr}
\toprule
Label & What the network reads & Test (\%) & Top-5 (\%) & Parameters & Training per seed (s) \\
\midrule
\multicolumn{6}{l}{\emph{Pixels: ResNet-18 (ImageNet) fine-tuned, then the small transformer}} \\
S3 & whole 224$\times$224 frame                  & 24.3 $\pm$ 1.9 & 54.0 & 11.4M & 2058 \\
S5 & two hand crops                                & 40.3 $\pm$ 2.1 & 68.7 & 11.5M & 942 \\
X1 & two hand crops + wrist location              & 41.7 $\pm$ 1.2 & 71.7 & 11.5M & 958 \\
\multicolumn{6}{l}{\emph{Pixels: video network}} \\
I3D & whole frames, 16 (Kinetics, fine-tuned)      & 26.3 $\pm$ 0.5 & 58.7 & 12.4M & 1371 \\
\midrule
\multicolumn{6}{l}{\emph{Keypoints: the small transformer, trained from scratch}} \\
S1 & hand keypoints + wrist location               & 59.0 $\pm$ 2.9 & 83.7 & 0.16M & 20 \\
X2 & hand + upper-body keypoints                   & 65.0 $\pm$ 1.6 & 85.0 & 0.17M & 23 \\
\bottomrule
\end{tabular}
\caption{Main results on the 100 test videos (mean $\pm$ seed standard deviation over 3 seeds). All rows use the same 748
training and 165 validation videos and the same 16 frames. Training time is on one RTX 3060 and excludes keypoint
extraction. A random guess scores 1.}
\label{tab:main}
\end{table*}

\begin{figure}[t]
\centering
\fig{figs/ladder_bars.png}{0.94\columnwidth}
\caption{Accuracy of each ladder step on the 100 test videos, with the 95\% bootstrap interval for one model (it is wide:
about $\pm$9 points). Compare models with the paired intervals in \cref{tab:paired}, which are narrower.}
\label{fig:bars}
\end{figure}

\begin{table*}[t]
\centering\small\setlength{\tabcolsep}{4pt}
\begin{tabular}{lrr}
\toprule
Comparison (A $-$ B) & Difference (points) & 95\% interval \\
\midrule
\textbf{Keypoint transformer, hands (matched)} $-$ ResNet, hand crops + wrist location, fine-tuned & $+17.3$ & $+9.0$ to $+25.3$ \\
\textbf{ResNet, hand crops + wrist location, fine-tuned} $-$ ResNet, hand crops, fine-tuned & $+1.3$ & $-3.7$ to $+6.3$ \\
\textbf{Keypoint transformer, hands + upper body} $-$ Keypoint transformer, hands (matched) & $+6.0$ & $+1.0$ to $+11.3$ \\
\midrule
Keypoint transformer, hands (matched) $-$ ResNet, hand crops, fine-tuned & $+18.7$ & $+9.7$ to $+27.3$ \\
Keypoint transformer, hands (matched) $-$ ResNet, whole frame, fine-tuned & $+34.7$ & $+25.0$ to $+44.0$ \\
ResNet, hand crops, fine-tuned $-$ ResNet, whole frame, fine-tuned & $+16.0$ & $+9.0$ to $+23.3$ \\
Keypoint transformer, hands (matched) $-$ Keypoint transformer, no wrist location & $+2.0$ & $-3.0$ to $+7.0$ \\
Keypoint transformer, tracked full-resolution $-$ Keypoint transformer, hands (matched) & $+8.0$ & $+2.0$ to $+14.3$ \\
Keypoint transformer, no depth $-$ Keypoint transformer, hands (matched) & $+3.3$ & $-1.0$ to $+7.7$ \\
Keypoint transformer, hands (matched) $-$ I3D, whole frame, fine-tuned & $+32.7$ & $+22.3$ to $+43.0$ \\
Keypoint transformer, hands + upper body $-$ ResNet, hand crops + wrist location, fine-tuned & $+23.3$ & $+15.3$ to $+31.3$ \\
ResNet, hand crops (black if missed), frozen $-$ ResNet, whole frame, frozen & $+8.3$ & $-0.7$ to $+17.3$ \\
ResNet, hand crops (black if missed), frozen $-$ ResNet, hand crops (borrowed box), frozen & $+8.3$ & $+2.0$ to $+15.0$ \\
ResNet, whole frame, fine-tuned $-$ ResNet, whole frame, frozen & $+8.7$ & $+2.0$ to $+15.3$ \\
ResNet, hand crops, fine-tuned $-$ ResNet, hand crops (borrowed box), frozen & $+24.7$ & $+17.0$ to $+32.3$ \\

\bottomrule
\end{tabular}
\caption{Paired differences in accuracy on the same 100 test videos, with 95\% bootstrap intervals over videos. Bold: the
three primary comparisons, fixed before the new runs finished. An interval that contains 0 means the difference is not
detectable at this sample size; it does not show that the models are equal.}
\label{tab:paired}
\end{table*}

\cref{tab:main} and \cref{fig:bars} give the ladder. On the same 16 frames, fine-tuned pixel models are clearly behind
the keypoint transformer. Cutting out the hands helps a ResNet by 16.0 points (95\% interval 9.0 to 23.3) over the whole
frame, so \emph{where} the network looks matters a lot. The keypoint model S1 is 18.7 points ahead of the crop model S5 (9.7 to 27.3), but S1 also knows where the hands are
in the frame and S5 does not. The primary comparison, keypoints against crops that also know the hand position (S1 vs.\ X1),
removes that difference: S1 is still ahead of X1 by 17.3 points (95\% interval 9.0 to 25.3), the first primary
comparison. Giving the crop model the wrist location changes its accuracy by only 1.3 points ($-3.7$ to $+6.3$), the second
primary comparison, so at this sample size location does not explain the gap. Adding the upper body lifts the keypoint model by 6.0 points (1.0 to 11.3), so
the hands alone leave information unused. The video network I3D, which reads the whole frame, scores 26.3 (top-5 58.7),
close to the whole-frame ResNet (24.3) and well below the keypoint model; keypoints are ahead of I3D by 32.7 points (22.3 to 43.0).
The keypoint model with the upper body is 23.3 points ahead of the crop model with location (15.3 to 31.3). \textbf{What these numbers do not show}: MediaPipe was trained for hands and
people, while the ResNet and I3D weights come from everyday photos and human actions, so the comparison is between pose
\emph{pipelines} and pixel \emph{pipelines}, not between representations in general.

\subsection{Cost}
Training the keypoint transformer takes about 20 seconds per seed, against 958 s for the crop ResNet with location, 1371 s
for I3D and 2058 s for the whole-frame ResNet (one RTX 3060, keypoint extraction excluded). Training is not the whole
cost, because keypoints must be extracted from every video. \cref{tab:speed} times everything from the video file to the
answer, one video at a time with nothing cached: reading and resizing the frames, MediaPipe where the pipeline needs it,
cropping, and the network. MediaPipe runs on the CPU (AMD Ryzen 5 5500), the networks on the GPU (RTX 3060); the networks
are untrained copies of the same architecture, since speed does not depend on the weights.

\begin{table}[t]
\centering\small\setlength{\tabcolsep}{2.5pt}
\begin{tabular}{lccccc}
\toprule
Pipeline & Test & Decode & MediaPipe & Network & Total \\
\midrule
Whole frame, ResNet (S3)           & 24.3 & 82 & --   & 21 & 102 \\
I3D, whole frames                  & 26.3 & 82 & --   & 19 & 101 \\
Hand crops + location, ResNet (X1) & 41.7 & 82 & 541  & 12 & 635 \\
Keypoints, hands (S1)              & 59.0 & 82 & 541  & 1  & 624 \\
Keypoints, hands + body (X2)       & 65.0 & 82 & 1272 & 1  & 1356 \\
Keypoints, tracked, full resolution & 67.0 & -- & 2122 & 1  & 2123 \\
\bottomrule
\end{tabular}
\caption{End-to-end time per video in milliseconds (mean over 99 test videos; the first video is a warm-up). ``MediaPipe''
for X2 is Hands (541) plus Pose (731); for the tracked row it is one pass over every frame of the full-resolution video,
which includes decoding. ``Network'' includes cropping.}
\label{tab:speed}
\end{table}

MediaPipe, not the network, is the cost: the keypoint transformer needs 1.3 ms, against 10.5 ms for the crop ResNet and
20.5 ms for the whole-frame ResNet, but extracting the keypoints takes 541 ms. The keypoint pipeline is therefore about six
times slower end to end than the whole-frame ResNet (624 vs.\ 102 ms) while scoring 35 points higher, and the crop ResNet
costs the same as the keypoint pipeline (635 vs.\ 624 ms) because its crops also need MediaPipe, while scoring 17 points
lower. The upper body doubles the cost (1356 ms), and tracking over the full-resolution video more than triples it
(2123 ms). These times are for a Python CPU implementation of MediaPipe at batch size one; an optimised or GPU deployment
would change the absolute numbers (\cref{fig:time}).

\begin{figure}[t]
\centering
\fig{figs/acc_vs_time.png}{0.94\columnwidth}
\caption{Test accuracy against end-to-end time per video (log scale), from \cref{tab:speed}. Keypoint pipelines (blue)
are more accurate but pay for MediaPipe; the pixel pipelines that need no detector (orange, left) are fast and less
accurate.}
\label{fig:time}
\end{figure}

\subsection{What the models see}
\cref{fig:examples} shows one test video as each group sees it.

\begin{figure}[t]
\centering
\fig{figs/examples.png}{0.86\columnwidth}
\caption{One test video, frames 2, 6, 10 and 14 of the 16. Top: the shared frame. Middle: the two hand crops (left hand
above right hand). Bottom: the keypoints found in that frame (hands in orange and green, upper body in grey).}
\label{fig:examples}
\end{figure}

\subsection{Where the models differ}
Counting a video as correct if a model gets it right in at least two of its three seeds, the crop ResNet with location (X1)
is right on 41 of the 100 test videos and the keypoint model (S1) on 60. They share 36 of these: the keypoint model gets 24
videos the crop model misses, the crop model gets only 5 that the keypoint model misses, and 35 videos neither gets right. So
the crop model's correct videos are nearly a subset of the keypoint model's. The pattern is the same against the whole-frame
ResNet (44 videos only the keypoint model gets, 4 only the ResNet) and I3D (43 and 5), and the keypoint model with the upper
body gets 29 videos that the crop model misses, against 3 the other way. The 100 videos cover 72 signs with at most three
videos each, so we do not report per-sign results.

%=====================================================================
\section{Ablations}\label{sec:ablations}
%=====================================================================
Ablations from the same matched setup are in \cref{tab:kp-abl,tab:pix-abl,tab:shuffle}; results from the earlier,
different protocol are kept apart in \cref{tab:earlier}. Two further tables (the graph network's twelve settings and the
design choices of the earlier keypoint models) are in the appendix (\cref{tab:graph,tab:design}).

\subsection{What the keypoint model is given}
\begin{table}[t]
\centering\small\setlength{\tabcolsep}{3pt}
\begin{tabular}{lcccc}
\toprule
Variant & Test (\%) & Top-5 & $\Delta$ vs.\ S1 & 95\% interval \\
\midrule
S1: hands + wrist location (base)        & 59.0 $\pm$ 2.9 & 83.7 & -- & -- \\
no wrist location                         & 57.0 $\pm$ 2.9 & 81.0 & $-2.0$ & $-7.0$ to $+3.0$ \\
no depth coordinate                       & 62.3 $\pm$ 0.9 & 87.0 & $+3.3$ & $-1.0$ to $+7.7$ \\
+ upper-body points (X2)                   & 65.0 $\pm$ 1.6 & 85.0 & $+6.0$ & $+1.0$ to $+11.3$ \\
tracked, full-resolution keypoints        & 67.0 $\pm$ 2.4 & 88.3 & $+8.0$ & $+2.0$ to $+14.3$ \\
graph network instead of transformer     & 26.7 $\pm$ 1.2 & 51.0 & -- & -- \\
\bottomrule
\end{tabular}
\caption{Keypoint-side ablations (16 frames, 3 seeds). The last two keypoint rows change the extraction, not the model:
tracking hands through the full-resolution video lets MediaPipe use neighbouring frames and more pixels. The graph
network follows Pose-TGCN \citep{li2020wlasl} on the 42 hand points with our own width and depth (best of six settings
by validation accuracy).}
\label{tab:kp-abl}
\end{table}

Only the upper body gives a detectable gain among the inputs of the matched keypoint model. Wrist location and depth
change the score by less than the interval: with 100 test videos we cannot tell whether they help. Tracking and full
resolution are worth about 8 points to the keypoint side, so a comparison that gives the skeleton full-resolution
tracked video and the images a shrunken frame overstates the skeleton's lead by about that much; this is why the
matched protocol was needed. Using 32 instead of 16 frames does not change the tracked-keypoint transformer (66.7 $\pm$ 1.9 with 32 frames in the
earlier run, 67.0 $\pm$ 2.4 with 16; same configuration, so the 16 frames used here lose nothing). The graph network is much
weaker than the transformer on the same keypoints (26.7 vs.\ 59.0). Across its twelve settings (\cref{tab:graph}), our
normalised features beat the raw coordinates the original used at every depth, tracked keypoints beat matched ones by
11 to 17 points (the transformer gained 8), more blocks do not help consistently, and several settings are unstable
across seeds (standard deviations up to 11 points).

\subsection{Choices in the picture models}
\begin{table}[t]
\centering\small\setlength{\tabcolsep}{3pt}
\begin{tabular}{lcc}
\toprule
Variant (ResNet-18) & Test (\%) & Top-5 \\
\midrule
whole frame, network frozen                       & 15.7 $\pm$ 1.7 & 30.7 \\
whole frame, fine-tuned (S3)                      & 24.3 $\pm$ 1.9 & 54.0 \\
hand crops, frozen, box borrowed when missed      & 15.7 $\pm$ 0.5 & 39.0 \\
hand crops, frozen, black when missed             & 24.0 $\pm$ 5.1 & 48.0 \\
hand crops, fine-tuned (S5)                        & 40.3 $\pm$ 2.1 & 68.7 \\
hand crops + wrist location, fine-tuned (X1)       & 41.7 $\pm$ 1.2 & 71.7 \\
\bottomrule
\end{tabular}
\caption{Picture-side ablations (16 frames, 3 seeds). ``Frozen'' keeps the ImageNet weights and trains only the
transformer on top. When MediaPipe finds no hand, the crop either reuses the box from the nearest frame where the hand
was found, or is black.}
\label{tab:pix-abl}
\end{table}

Fine-tuning helps the whole-frame ResNet by 8.7 points (2.0 to 15.3) and the crop ResNet by 24.7 (17.0 to 32.3), and
hand crops beat the whole frame once the network is fine-tuned (+16.0, 9.0 to 23.3). Replacing a borrowed box by a black
crop raises the frozen ResNet from 15.7 to 24.0 (+8.3, 2.0 to 15.0): a borrowed box shows whatever lies where the hand
used to be. Adding wrist location to the fine-tuned crops (X1) changes the score by 1.3 points ($-3.7$ to $+6.3$), not detectable;
for the keypoint model the same information was also not detectable ($-2.0$ when removed).

\subsection{Is movement used?}
\begin{table}[t]
\centering\small\setlength{\tabcolsep}{3pt}
\begin{tabular}{lccc}
\toprule
Model & Normal order & Shuffled & Change \\
\midrule
S1: keypoints, hands + location     & 59.0 & 62.1 & $+3.1$ \\
keypoints, no wrist location        & 57.0 & 56.8 & $-0.2$ \\
X2: keypoints, hands + upper body   & 65.0 & 65.1 & $+0.1$ \\
X1: hand crops + location            & 41.7 & 39.9 & $-1.8$ \\
\bottomrule
\end{tabular}
\caption{Test accuracy (\%) with the 16 frames of every test video in their normal order or shuffled (mean over three fixed
shuffles, the same for all models). The model is trained on ordered videos and is not retrained.}
\label{tab:shuffle}
\end{table}

Shuffling the frames does not lower the keypoint models' accuracy (\cref{tab:shuffle}), so as trained they depend on
which hand poses occur, not on their order. The crop ResNet loses 1.8 points, which is within the noise of 100 videos (we
did not compute an interval for it). This is a one-sided check: shuffled videos differ from the training data, so a
drop would have been hard to interpret, and no drop does not show that order carries no information for a model that
is trained to use it.

\subsection{Earlier protocol and other model types}
\begin{table}[t]
\centering\small\setlength{\tabcolsep}{3pt}
\begin{tabular}{lcc}
\toprule
Model (earlier protocol) & Test (\%) & Val (\%) \\
\midrule
\multicolumn{3}{l}{\emph{Keypoints from tracked full-resolution video, 32 frames}} \\
XGBoost, best on validation                         & 39.3 $\pm$ 1.2 & 47.5 \\
MLP, best on validation$^\ast$                       & 60.0 & 63.6 \\
GRU, best on validation$^\ast$                       & 60.7 & 66.3 \\
LSTM, best on validation$^\ast$                      & 63.3 & 63.4 \\
transformer, best on validation (mean over frames)    & 66.7 $\pm$ 1.9 & 72.3 \\
\quad base (one layer, last token, no augmentation)   & 68.0 $\pm$ 1.4 & 70.7 \\
\quad base with augmentation                          & 70.0 $\pm$ 2.2 & 71.9 \\
\quad two layers                                      & 67.0 & 70.1 \\
\quad two layers, no wrist location                   & 62.7 & 67.3 \\
\midrule
\multicolumn{3}{l}{\emph{Pixels, 112$\times$112, 16 frames, hands found in the full-resolution video}} \\
ResNet-18, hand crops, fine-tuned                    & 47.0 $\pm$ 3.7 & 46.5 \\
ResNet-18, whole frame, fine-tuned                   & 15.7 $\pm$ 1.2 & 24.6 \\
R(2+1)D \citep{tran2018closer}, whole frame, fine-tuned & 38.7 $\pm$ 3.1 & 37.4 \\
R(2+1)D, one crop per video, fine-tuned              & 34.3 $\pm$ 3.4 & 42.8 \\
\bottomrule
\end{tabular}
\caption{Results from the earlier protocol, in which the keypoints came from full-resolution video (with tracking) and the
pixel models read 112$\times$112 frames or crops. They are \emph{not} comparable with \cref{tab:main}.
$^\ast$Narelle's models, each the mean over three seeds (seed standard deviation up to 5.0); her per-video
predictions are not in this repository, so they are not in the paired comparisons. The R(2+1)D crop row fine-tuned longer with three learning rates
($3\times10^{-5}$, $10^{-4}$, $3\times10^{-4}$) gave 27, 46 and 36 on the test videos in one seed each, so tuning
changes a video network by more than the gaps between most rows here.}
\label{tab:earlier}
\end{table}

Under the earlier protocol, mean pooling or attention over time was what mattered most for keypoint models; once a
model aggregates over time, the transformer, GRU, LSTM and MLP are within noise of each other, and XGBoost on the
flattened sequence is the weakest (\cref{tab:earlier}). Hand crops cut from full-resolution video help a fine-tuned
ResNet more than crops cut from the 224-pixel frame (47.0 vs.\ 40.3, in different protocols and so not a tested
difference), and the strongest video-network result needed a learning-rate search, which we did not apply to the
ResNet or I3D rows of \cref{tab:main}.

\cref{tab:design} (appendix) lists the design choices of the earlier keypoint models. \textbf{Augmentation} helps models
that read a flat vector most (+8 to +17 points for the linear model, MLP and XGBoost) and the transformer least (+2).
\textbf{Temporal aggregation} is the largest gain for the MLP (+16), linear model (+17), LSTM (+12) and GRU (+7), but costs
XGBoost 4 points and the transformer 3 (for the transformer, validation accuracy rises while test accuracy falls, so we do
not read it as a loss).
\textbf{Wrist location} adds 2 to 6 points for every model that was tested with and without it. For the pixel models of
that protocol, the frozen R(2+1)D scores 10.7 to 12.3 whatever the number of frames (8 or 16) or the input (whole frame
or one crop), and a small ConvNet trained from scratch stays between 1.0 and 1.7, which is chance (1): without pretraining the
pixel models learn nothing.

\subsection{Limits that bear on these results}\label{sec:limits}
One recipe per group and no search; 16 frames of 224$\times$224 (hands about 24 pixels wide); hands (and in X2
the upper body) but no face; a 100-video test set that contains 72 of the 100 signs, with 97 of its videos by signers
seen in training; half of WLASL100; MediaPipe is itself trained on hands, the pixel networks are not. One test video (24960, ``go'') is an exact copy of a training video (24961): its frames are pixel-identical. Comparing every
test video with every training and validation video by frame content found no other pair with the same sign. We kept the
official test set, so this can raise a score by at most one point. Narelle's MLP, GRU and LSTM were trained with an
earlier recipe, and their per-video predictions are not available here for paired comparisons.

\bibliography{refs}
\bibliographystyle{icml2025}

\appendix
\section{Additional ablation tables}

\begin{table}[h]
\centering\footnotesize\setlength{\tabcolsep}{2.5pt}
\begin{tabular}{llccc}
\toprule
Keypoints & Input & 2 blocks & 4 blocks & 8 blocks \\
\midrule
matched  & raw coordinates & 8.3 $\pm$ 3.9 (13.5)  & 18.3 $\pm$ 1.2 (20.2) & 14.0 $\pm$ 1.4 (19.8) \\
matched  & normalised      & \textbf{26.7 $\pm$ 1.2 (30.5)} & 22.3 $\pm$ 11.0 (30.3) & 22.7 $\pm$ 7.4 (29.3) \\
tracked  & raw coordinates & 25.3 $\pm$ 4.6 (31.7) & 30.7 $\pm$ 0.5 (36.0) & 31.3 $\pm$ 0.5 (40.0) \\
tracked  & normalised      & 37.3 $\pm$ 2.9 (40.6) & 36.3 $\pm$ 0.5 (43.4) & \textbf{39.3 $\pm$ 4.0 (45.7)} \\
\bottomrule
\end{tabular}
\caption{The graph network (Pose-TGCN style) in its twelve settings: test accuracy (\%) $\pm$ seed standard deviation, with
validation accuracy in brackets. Bold: the setting chosen on validation accuracy for each kind of keypoints. ``Matched''
keypoints come from the shared 224-pixel frames; ``tracked'' from MediaPipe tracking through the full-resolution video. Raw
coordinates are the 2D image coordinates, as in the original; ``normalised'' are the features of the transformer. Parameters:
0.16M, 0.27M and 0.50M for 2, 4 and 8 blocks (0.16M, 0.28M, 0.51M normalised).}
\label{tab:graph}
\end{table}

\begin{table}[h]
\centering\footnotesize\setlength{\tabcolsep}{3pt}
\begin{tabular}{lccc}
\toprule
Model (earlier protocol) & Augmentation & Temporal aggregation & Wrist location \\
\midrule
Linear$^\ast$      & 18.7 $\to$ 36.0 & 36.0 $\to$ 52.7 & 16.0 $\to$ 18.7 \\
MLP$^\ast$         & 30.3 $\to$ 44.0 & 44.0 $\to$ 60.0 & 28.3 $\to$ 30.3 \\
XGBoost            & 37.3 $\to$ 45.7 & 45.7 $\to$ 41.3 & 31.3 $\to$ 37.3 \\
GRU$^\ast$         & 43.0 $\to$ 53.3 & 53.3 $\to$ 60.7 & -- \\
LSTM$^\ast$        & 47.3 $\to$ 50.0 & 51.0 $\to$ 63.3 & -- \\
Transformer        & 68.0 $\to$ 70.0 & 70.0 $\to$ 66.7 & 62.7 $\to$ 67.0 \\
\bottomrule
\end{tabular}
\caption{Design choices of the earlier keypoint models (tracked full-resolution keypoints, 32 frames): test accuracy (\%)
before $\to$ after adding the choice, one choice at a time, mean over three seeds. Augmentation: random rotation, zoom,
shift and trimming of the clip. Temporal aggregation: mean and standard deviation over time (linear, MLP, XGBoost), a mean
over time instead of the last step (GRU, LSTM, transformer). Wrist location: the four wrist-position numbers (for the
transformer, with two layers). $^\ast$Narelle's models.}
\label{tab:design}
\end{table}

\end{document}
